\scanpage{640}
\mspage{28}

dans $M$ (laquelle sera forcément unique), il f. et
s. que $M$ et $W$ satisfassent les conditions
(2.8.1 a) b)), et que de plus on ait

\[
\left\{\begin{aligned}
\text{a)}\quad &TC\overset{\mathrm{def}}{=}\mathrm{Mono}\cap W\ \text{est stable par cochg}^{\mathrm t}\text{ de base}\\
\text{b)}\quad &TF\overset{\mathrm{def}}{=}(\mathrm{Mono})_*\subset W,\quad\text{i.e. }W\text{ contient}\\
 &\text{les flèches de type injectif dans}\\
 &M.
\end{aligned}\right.\tag*{(2.8.2)}
\]
\deleted{\unsure{(2.8.3)}}

% slanted
\margin{On dit \unsure{aussi|ainsi} que le localiseur $W$ est très spécial, s'il définit une str. de Quillen très spéciale.}

On trouve maintenant le

\uline{\deleted{\unsure{Corollaire}}Théorème} 2.9. Soit $(W,C\unsure{=|\underset{\vert}{=}}\mathrm{Mono},F=(TC)_*=(W\cap\mathrm{Mono})_*)$
un triple de Quillen très spécial dans une
catégorie $M$ (2.8.). Alors \deleted{pour}

$1^\circ$) Pour \unsure{tout} $I$ dans Cat, $(W(I),\deleted{\unsure{W}}C(I),F[I])$
est un \deleted{couple}\inserted{triple} de Quillen très spécial
dans $M(I)$, et on a \deleted{\unsure{si $v\colon I_0\unsure{-|\longrightarrow}I$ est} défini \unsure{comme}}\inserted{en particulier}

\[
 C(I)=\operatorname{Mono}(M(I)),\tag*{(2.9.1)}
\]

et de plus

\[
 F[I]\cap W(I)=(TF)[I].\tag*{(2.9.2)}
\]

$2^\circ$) Pour toute flèche $u\colon I\unsure{-|\longrightarrow}I'$ dans Cat,
et \deleted{\unsure{où|au}} \deleted{\unsure{le}}le couple de foncteurs adjoints

\diagram{\begin{tikzcd}[ampersand replacement=\&]
\text{(2.9.3)} \& M(I) \arrow[r,shift right=0.8ex,"u_*"'] \& M(I') \arrow[l,shift right=0.8ex,"u^*"'] ,
\end{tikzcd}}

on a

\[
\left\{\begin{aligned}
 u^*(W(I'))&\subset W(I)\\
 u^*(C(I'))&\subset C(I)\\
 u^*(TC(I'))&\subset TC(I)\\
 u_*(F[I])&\subset F[I']\\
 u_*(TF[I])&\subset TF[I']
\end{aligned}\right.,\tag*{(2.9.\unsure{4|3})}
\]

enfin $u_*$ induit un foncteur
% archive sheet number: 639
